\section{AlphaCode 2：更强基础模型加学习式选择器}

\subsection{从自建代码模型切换到 Gemini Pro}

\videofigure{figures/alphacode2-overview.jpg}{AlphaCode 2 用 Gemini Pro 替换原始代码模型，并在采样与筛选端加入更多学习组件。}{00:25:14--00:26:48}

AlphaCode 2 不再从头预训练代码模型，而是利用更强的通用 foundation model，再针对竞赛编程微调多个 solver family。这将“知识与代码能力”交给基础模型，把系统创新集中在多样采样与排序。

\videofigure{figures/alphacode2-contributions.jpg}{AlphaCode 2 的主要变化包括 Gemini 微调、多个 solver family，以及提交前的正确性 scoring model。}{00:26:48--00:27:42}

\begin{importantbox}{基础模型与搜索器是乘法关系}
更强模型提高单次命中率 $p$；更好的搜索和选择提高预算利用率。二者共同决定最终 solve rate，不能把全部进步归因于“模型更大”或“采样更多”。
\end{importantbox}

\subsection{Scoring model 学习“哪个候选更可能正确”}

\videofigure{figures/scoring-model.jpg}{独立 scoring model 根据题目与候选程序估计正确性，用于提交前排序。}{00:27:42--00:29:18}

设候选为 $y_i$，打分器输出 $s_\phi(x,y_i)$。最终系统不是按 solver 的生成概率排序，而是组合：

$$
\operatorname{rank}(y_i)=
\lambda s_\phi(x,y_i)+(1-\lambda)d(y_i,\mathcal S),
$$

其中 $d(y_i,\mathcal S)$ 表示相对已选集合的语义多样性。高分候选更可能正确，多样性项避免全部提交来自同一错误簇。

\begin{warningbox}{生成概率不能直接当正确性概率}
模型偏好的往往是常见、流畅、模板化的代码，而不是通过所有边界测试的代码。课程结尾进一步指出，语言模型常对错误答案过度自信，概率校准本身就是开放问题。
\end{warningbox}

\videofigure{figures/alphacode2-sampling.jpg}{AlphaCode 2 聚合多个 solver family 的大规模样本，经测试、聚类和打分后保留少量提交。}{00:29:18--00:30:42}

\subsection{用更少样本达到更高 solve rate}

\videofigure{figures/solve-rate-scaling.jpg}{AlphaCode 2 的 solve rate 随采样预算增长，并以远少于旧系统的样本达到可比或更好结果。}{00:30:42--00:32:10}

AlphaCode 2 的重要工程结论是：提高 proposal quality 与 selector quality 后，不必永远把样本数推到一百万。对简单题，较小预算已能覆盖正确解；对难题，额外采样才更有价值。

\videofigure{figures/alphacode2-percentile.jpg}{AlphaCode 2 在模拟竞赛评估中达到约第 85 百分位，明显超过原始 AlphaCode。}{00:37:18--00:39:02}

\videofigure{figures/alphacode2-tradeoffs.jpg}{AlphaCode 2 提升样本效率和候选质量，但运行成本高、方法仍高度依赖代码可执行验证。}{00:39:02--00:40:02}

\begin{knowledgebox}{Adaptive compute 是自然下一步}
先用轻量模型估计题目难度和当前候选置信度，再动态决定采样数、温度、solver family 与是否进入串行修正，可把平均成本从“每题固定一百万样本”降到按需分配。
\end{knowledgebox}

\subsection{本章小结}

AlphaCode 2 通过更强基础模型、多 solver family 和 learned scorer 改善 proposal 与 selection，说明测试时计算的价值不只来自规模，还来自任务自适应和更准确的候选排序。

